% from ega2-1-fr.tex
\oldpage[II]{56}

\begin{corollary}[3.2.8]
\label{II.3.2.8-fr}
Supposons $\mathcal{S}$ engendrée par $\mathcal{S}_1$. Pour tout
$\mathcal{S}$-Module gradué $\mathcal{M}$ et tout $n\in\mathbf{Z}$, on a
\[
\label{II.3.2.8.1-fr}
  (\mathcal{M}(n))^{\sim}=\widetilde{\mathcal{M}}(n)
\tag{3.2.8.1}
\]
à un isomorphisme canonique près.
\end{corollary}

\begin{proof}
Cela résulte des propriétés correspondantes pour $Y$ affine
(\hyperref[II.2.5.14-fr]{2.5.14} et
\hyperref[II.2.5.15-fr]{2.5.15}) ainsi que de
(\hyperref[II.2.8.11-fr]{2.8.11}).
\end{proof}

\par\medskip\noindent\textit{Remarques (3.2.9). ---\ }
\label{II.3.2.9-fr}
\begin{enumerate}
  \item[\rm(i)] Si $\mathcal{S}=\mathcal{A}[T]$, $\mathcal{A}$ étant une
    $\mathcal{O}_Y$-Algèbre quasi-cohérente
    (\hyperref[II.3.1.7-fr]{3.1.7}), on vérifie immédiatement que tous les
    $\mathcal{O}_X$-Modules inversibles $\mathcal{O}_X(n)$ sont canoniquement
    isomorphes à $\mathcal{O}_X$.

    En outre, soit $\mathcal{N}$ un $\mathcal{A}$-Module quasi-cohérent, et
    posons $\mathcal{M}=\mathcal{N}\otimes_{\mathcal{A}}\mathcal{A}[T]$. Il
    résulte alors de (\hyperref[II.3.2.3-fr]{3.2.3}) et de
    (\hyperref[II.3.1.7-fr]{3.1.7}) que dans l'identification canonique de
    $X=\operatorname{Proj}(\mathcal{A}[T])$ et de
    $X'=\operatorname{Spec}(\mathcal{A})$, le $\mathcal{O}_X$-Module
    $\widetilde{\mathcal{M}}$ s'identifie au $\mathcal{O}_{X'}$-Module
    $\widetilde{\mathcal{N}}$ associé à $\mathcal{N}$ (au sens de
    (\hyperref[II.1.4.3-fr]{1.4.3})).
  \item[\rm(ii)] Soit $\mathcal{S}$ une $\mathcal{O}_Y$-Algèbre graduée
    quelconque, $\mathcal{S}'$ la $\mathcal{O}_Y$-Algèbre graduée telle que
    $\mathcal{S}'_0=\mathcal{O}_Y$, $\mathcal{S}'_n=\mathcal{S}_n$ pour tout
    $n>0$~; l'isomorphisme canonique de
    $X=\operatorname{Proj}(\mathcal{S})$ sur
    $X'=\operatorname{Proj}(\mathcal{S}')$
    (\hyperref[II.3.1.8-fr]{3.1.8, (ii)}) identifie $\mathcal{O}_X(n)$ et
    $\mathcal{O}_{X'}(n)$ pour tout $n\in\mathbf{Z}$~: cela résulte de la même
    proposition pour le cas affine (\hyperref[II.2.5.16-fr]{2.5.16}) et du
    fait que ces identifications, pour les ouverts affines de $Y$, commutent
    avec les opérations de restriction. De même, soit
    $X^{(d)}=\operatorname{Proj}(\mathcal{S}^{(d)})$~; l'isomorphisme
    canonique de $X$ sur $X^{(d)}$
    (\hyperref[II.3.1.8-fr]{3.1.8, (i)}) identifie $\mathcal{O}_X(nd)$ et
    $\mathcal{O}_{X^{(d)}}(n)$ pour tout $n\in\mathbf{Z}$.
\end{enumerate}

\begin{proposition}[3.2.10]
\label{II.3.2.10-fr}
Soient $\mathcal{L}$ un $\mathcal{O}_Y$-Module inversible, $g$ l'isomorphisme
canonique
$X_{(\mathcal{L})}=\operatorname{Proj}(\mathcal{S}_{(\mathcal{L})})
\to X=\operatorname{Proj}(\mathcal{S})$
(\hyperref[II.3.1.8-fr]{3.1.8, (iii)}). Pour tout entier $n\in\mathbf{Z}$,
$g_*(\mathcal{O}_{X_{(\mathcal{L})}}(n))$ est canoniquement isomorphe à
$\mathcal{O}_X(n)\otimes_Y\mathcal{L}^{\otimes n}$.
\end{proposition}

\begin{proof}
Supposons d'abord $Y$ affine, d'anneau $A$, et
$\mathcal{L}=\widetilde{L}$, où $L$ est un $A$-module libre monogène. Avec les
notations de la démonstration de
(\hyperref[II.3.1.8-fr]{3.1.8, (iii)}) on définit, pour $f\in S_d$, un
isomorphisme de $(S(n))_{(f)}\otimes_A L^{\otimes n}$ sur
$(S_{(L)}(n))_{(f\otimes c^{\otimes d})}$ en faisant correspondre à
$(x/f^k)\otimes c^{\otimes n}$, où $x\in S_{kd+n}$, l'élément
$(x\otimes c^{\otimes(n+kd)})/(f\otimes c^{\otimes d})^k$~; il est immédiat
que cet isomorphisme est indépendant du générateur $c$ choisi pour $L$~; en
outre, les isomorphismes ainsi définis pour chaque $f\in S_+$ sont compatibles
avec les opérations de restriction $D_+(f)\to D_+(fg)$. Enfin, dans le cas
général, on voit sans peine en revenant aux définitions
(\hyperref[II.3.1.1-fr]{3.1.1}) que les isomorphismes ainsi définis pour
chaque ouvert affine $U$ de $Y$ sont compatibles avec le passage de $U$ à un
ouvert affine $U'\subset U$.
\end{proof}

\subsection{$\mathcal{S}$-Module gradué associé à un faisceau sur
$\operatorname{Proj}(\mathcal{S})$.}
\label{subsection:II.3.3-fr}

\emph{On suppose dans tout ce numéro que la $\mathcal{O}_Y$-Algèbre graduée
$\mathcal{S}$ est engendrée par $\mathcal{S}_1$
(\hyperref[II.3.1.9-fr]{3.1.9}).} Rappelons qu'en raison de
(\hyperref[II.3.1.8-fr]{3.1.8, (i)}), cette restriction n'est pas essentielle,
moyennant les conditions de finitude (\hyperref[II.3.1.10-fr]{3.1.10}).

\begin{env}[3.3.1]
\label{II.3.3.1-fr}
Soit $p$ le morphisme structural $X=\operatorname{Proj}(\mathcal{S})\to Y$.
Pour tout $\mathcal{O}_X$-Module $\mathcal{F}$, nous poserons
\[
\label{II.3.3.1.1-fr}
  \Gamma_*(\mathcal{F})
  =\bigoplus_{n\in\mathbf{Z}}p_*(\mathcal{F}(n))
\tag{3.3.1.1}
\]

